%% toc: P2 21 21.8
\section{Proofs for finite versus unrestricted validity}
\label{sec:annex-p02-w02-p02-036-finite-validity}

Unless an item says otherwise, \(\Lang\) is a finite relational signature,
\(X=\operatorname{Sent}(\Lang)\), and the classes
\(\mathrm{Fin\text{-}sat}\), \(\mathrm{Non\text{-}sat}\) and
\(\mathrm{Inf\text{-}axioms}\) are those of
Definition~\ref{def:p02-w02-p02-036-classes}.  Sets of sentences are sets of
codes under the fixed effective coding of \cref{sec:decision-hierarchy-01};
under that coding \(X\) is a decidable set.  Two facts about computably
enumerable sets are used repeatedly: a set is computably enumerable exactly
when some Turing machine recognizes it (\cref{thm:ce-equivalent-forms}), and a
set is decidable exactly when both it and its complement are computably
enumerable (the complement criterion in \cref{sec:turing-machines-03}).  A
union of two computably enumerable sets is computably enumerable: run
recognizers for the two sets alternately, one step at a time, and accept when
either accepts.

\begin{annexitem}{annex:p02-w02-p02-036-infinity-axiom}{thm:p02-w02-p02-036-infinity-axiom}{The sentence \(\prec_K\) has only infinite models}
\begin{proof}
\emph{A model.}  Read \(K\) as \(<\) on \(\mathbb N\).  Given \(x\), take
\(u=x+1\).  Then \(\neg x<x\) and \(x<x+1\) hold, and for every \(y\), if
\(y<x\) then \(y<x+1\).  So the quantifier-free part holds for every \(y\),
and \((\mathbb N,<)\entails\prec_K\).

\emph{No finite model.}  Let \(\mathfrak A\entails\prec_K\) and write \(K\) for
\(K^{\mathfrak A}\).  Choose \(a_0\) in the domain, which is nonempty.  Given
\(a_i\), the sentence \(\prec_K\) with \(x\mapsto a_i\) supplies an element
\(a_{i+1}\) such that \(Ka_ia_{i+1}\) holds and, for every element \(y\),
\(Kya_i\) implies \(Kya_{i+1}\).  This defines \(a_0,a_1,a_2,\ldots\).  We show
\(Ka_ia_j\) for all \(i<j\), by induction on \(j\).  For \(j=i+1\) it holds by
the choice of \(a_{i+1}\).  If \(Ka_ia_j\) holds, apply the property of
\(a_{j+1}\) with \(y=a_i\): from \(Ka_ia_j\) it gives \(Ka_ia_{j+1}\).  If
\(a_i=a_j\) for some \(i<j\), then \(Ka_ia_i\), which contradicts the conjunct
\(\neg Kxx\) at \(x\mapsto a_i\).  So the \(a_i\) are pairwise distinct, and
the domain is infinite.
\end{proof}
\end{annexitem}

\begin{annexitem}{annex:p02-w02-p02-036-finsat-ce}{thm:p02-w02-p02-036-finsat-ce}{Finite model checking and enumerability}
\begin{proof}
\emph{Item 1.}  Let \(\mathcal M\) have domain \(\{0,\ldots,n-1\}\).  For a
formula \(\theta\) and a finite map \(s_0\) from a finite set of variables
containing \(\FV(\theta)\) into \(\{0,\ldots,n-1\}\), compute a truth value
\(V(\theta,s_0)\) by recursion on the complexity of \(\theta\):
\(V(x=y,s_0)\) is true iff \(s_0(x)=s_0(y)\); \(V(R_i(x_1,\ldots,x_{m_i}),s_0)\)
is true iff the tuple \((s_0(x_1),\ldots,s_0(x_{m_i}))\) occurs in the given
list for \(R_i\); \(V(\bottom,s_0)\) is false; the connectives are evaluated by
their truth tables from the values of the immediate subformulas; and
\(V(\forall x\,\theta,s_0)\), respectively \(V(\exists x\,\theta,s_0)\), is true
iff \(V(\theta,s_0[x\mapsto a])\) is true for every, respectively for some,
\(a<n\).  Each case calls the procedure only on immediate subformulas, which
have smaller complexity, and only finitely many times, since there are \(n\)
values of \(a\).  So the recursion terminates.  Its clauses are the
satisfaction clauses of \cref{sec:fol-semantics-02} with the quantifiers
ranging over the finite domain.  Hence, by induction on complexity and by
\cref{lem:fol-coincidence}, \(V(\theta,s_0)\) is true iff
\(\mathcal M,s\entails\theta\) for every assignment \(s\) extending \(s_0\).
For a sentence \(\varphi\), take \(s_0\) empty.

\emph{Item 2.}  For each \(n\geq1\) the structures with domain
\(\{0,\ldots,n-1\}\) are given by choosing, for each \(i\leq r\), a subset of
the finite set \(\{0,\ldots,n-1\}^{m_i}\).  There are finitely many choices,
and they can be listed effectively.  A machine that, on input \(\varphi\),
first checks \(\varphi\in X\), then runs through \(n=1,2,\ldots\) and, for each
\(n\), through these finitely many structures, and accepts as soon as item~1
reports a model of \(\varphi\), accepts exactly the sentences with a model on
some \(\{0,\ldots,n-1\}\).  By
Convention~\ref{def:p02-w02-p02-036-finite-data} and \cref{lem:iso-invariance},
these are exactly the members of \(\mathrm{Fin\text{-}sat}(X)\).  Hence that set
is computably enumerable.

\emph{Item 3.}  A sentence \(\varphi\in X\) has no model iff \(\neg\varphi\)
is true in every \(\Lang\)-structure, that is, iff \(\neg\varphi\) is valid.
By \cref{lem:reduct-invariance}, truth of an \(\Lang\)-sentence in a
structure of a larger signature agrees with truth in its \(\Lang\)-reduct, and
every \(\Lang\)-structure has expansions, so validity of \(\neg\varphi\) does
not depend on the signature in which it is read.  The set of valid sentences is
computably enumerable (\cref{thm:fol-validity-enumerable}).  A machine that
checks \(\varphi\in X\) and then runs a recognizer for the valid sentences on
\(\neg\varphi\) recognizes \(\mathrm{Non\text{-}sat}(X)\).

\emph{Item 4.}  A sentence \(\varphi\in X\) fails to be finitely valid iff some
finite \(\Lang\)-structure satisfies \(\neg\varphi\), that is, iff
\(\neg\varphi\in\mathrm{Fin\text{-}sat}(X)\).  Running the recognizer of
item~2 on \(\neg\varphi\) after checking \(\varphi\in X\) recognizes the
complement in \(X\) of the set of finitely valid sentences.
\end{proof}
\end{annexitem}

\begin{annexitem}{annex:p02-w02-p02-036-run-sentence}{thm:p02-w02-p02-036-run-sentence}{Accepting runs and finite models}
\begin{proof}
Write \(w=a_0\cdots a_{k-1}\) and let the computation of \(\mathcal T\) on
\(w\) be \(c_0,c_1,\ldots\), with \(c_i=(q_i,T_i,h_i)\) and
\(c_0=(q_0,T_w,0)\) as in \cref{sec:turing-machines-01}.  Each step changes
the head position by exactly one, so \(\lvert h_i\rvert\leq i\).

\emph{Step 1: finite linear orders.}  Let \(<\) be a strict linear order of a
finite nonempty set \(A\), that is, a relation satisfying the order clause.
There is a least element: start anywhere and, while the current element has
a smaller one, move to a smaller one.  By transitivity every element visited
is below all earlier ones, and by irreflexivity no element is visited twice,
so the walk stops at an element with nothing below it, which is least by
linearity.  Removing the least element and repeating lists \(A\) as
\(e_0<e_1<\cdots<e_N\), and by transitivity \(e_i<e_j\) iff \(i<j\).  Then
\(\mathsf{Fst}(x)\) holds iff \(x=e_0\), and \(\mathsf{Nx}(x,y)\) holds iff
\(x=e_i\) and \(y=e_{i+1}\) for some \(i<N\): if \(x=e_i\), \(y=e_j\) and \(x<y\),
then \(i<j\), and some \(e_l\) lies strictly between them iff \(j>i+1\).

\emph{Step 2: a finite model gives an accepting run.}  Let \(\mathcal M\) be a
finite model of \(\psi_{\mathcal T,w}\), with domain listed as in Step~1.  By
the determinacy clause, for each \(i\leq N\) there are a unique state
\(\bar q_i\) with \(\mathsf{St}_{\bar q_i}(e_i)\), a unique index \(\bar h_i\) with
\(\mathsf{Hd}(e_i,e_{\bar h_i})\), and for each \(j\leq N\) a unique symbol
\(\bar T_i(j)\) with \(\mathsf{Cl}_{\bar T_i(j)}(e_i,e_j)\).  By the
acceptance clause some \(e_i\) satisfies \(\mathsf{St}_{q_{\mathrm{acc}}}\), so
there is a least \(i_0\leq N\) with
\(\bar q_{i_0}\in\{q_{\mathrm{acc}},q_{\mathrm{rej}}\}\).  Put \(o=\bar h_0\).

We show by induction on \(i\leq i_0\) that \(c_i\) exists and that
(a)~\(q_i=\bar q_i\); (b)~\(0\leq o+h_i\leq N\) and \(\bar h_i=o+h_i\); and
(c)~\(\bar T_i(j)=T_i(j-o)\) for every \(j\leq N\).

For \(i=0\): by Step~1, \(\mathsf{Fst}(e_0)\), so the start clause gives
\(\mathsf{St}_{q_0}(e_0)\), hence \(\bar q_0=q_0\).  Item~(b) holds since
\(h_0=0\).  The start clause applied to \(t=e_0\) and \(p=e_o\) gives
\(\mathsf{In}_w(e_0,e_o)\).  If \(k=0\), every cell holds \(\sqcup\) at
\(e_0\), and \(T_w(j-o)=\sqcup\) for every \(j\).  If \(k>0\), take witnesses
\(y_0,\ldots,y_{k-1}\).  Then \(y_0=e_o\), and from
\(\mathsf{Nx}(y_i,y_{i+1})\) and Step~1, induction on \(i\) gives
\(y_i=e_{o+i}\), so in particular \(o+k-1\leq N\).  The clause
\(\mathsf{Cl}_{a_i}(e_0,y_i)\) gives \(\bar T_0(o+i)=a_i=T_w(i)\).  For
\(j\notin\{o,\ldots,o+k-1\}\) the last conjunct gives
\(\bar T_0(j)=\sqcup\), and \(T_w(j-o)=\sqcup\) because \(j-o\) lies
outside \(\{0,\ldots,k-1\}\).  This is~(c).

Suppose the claims hold for some \(i<i_0\).  Since \(i<i_0\), the state
\(q_i=\bar q_i\) is not halting, so \(c_{i+1}\) exists.  Let
\(a=T_i(h_i)\) and \(\delta(q_i,a)=(q',b,d)\), so that
\(c_{i+1}=(q',T_{i+1},h_{i+1})\), where \(T_{i+1}(h_i)=b\), \(T_{i+1}\) agrees
with \(T_i\) elsewhere, and \(h_{i+1}=h_i+1\) if \(d=R\), \(h_{i+1}=h_i-1\) if
\(d=L\).  Put \(p=e_{o+h_i}\).  By (b) and (c),
\(\mathsf{Hd}(e_i,p)\) and \(\mathsf{Cl}_a(e_i,p)\), and
\(\mathsf{St}_{q_i}(e_i)\) holds by (a).  The step clause for \((q_i,a)\) gives
\(t'\) and \(p'\) with the listed properties.  From \(\mathsf{Nx}(e_i,t')\) and
Step~1, \(t'=e_{i+1}\), so \(i+1\leq N\).  From \(\mathsf{Mv}_d(p,p')\) and
Step~1, \(p'=e_{o+h_i+1}\) if \(d=R\) and \(p'=e_{o+h_i-1}\) if \(d=L\);
in both cases \(p'=e_{o+h_{i+1}}\), and \(0\leq o+h_{i+1}\leq N\).  The
conjunct \(\mathsf{St}_{q'}(e_{i+1})\) and uniqueness give
\(\bar q_{i+1}=q'=q_{i+1}\), which is~(a).  The conjunct
\(\mathsf{Hd}(e_{i+1},p')\) and uniqueness give
\(\bar h_{i+1}=o+h_{i+1}\), which is~(b).  The conjunct
\(\mathsf{Cl}_b(e_{i+1},p)\) gives \(\bar T_{i+1}(o+h_i)=b=T_{i+1}(h_i)\).
For \(j\neq o+h_i\), apply the frame conjunct with \(u=e_j\) and
\(c=\bar T_i(j)\): it gives \(\mathsf{Cl}_{\bar T_i(j)}(e_{i+1},e_j)\), so
\(\bar T_{i+1}(j)=\bar T_i(j)=T_i(j-o)=T_{i+1}(j-o)\), the last because
\(j-o\neq h_i\).  This is~(c).

At \(i=i_0\), item (a) gives \(q_{i_0}=\bar q_{i_0}\), a halting state.  The
acceptance clause says that no element satisfies
\(\mathsf{St}_{q_{\mathrm{rej}}}\), so \(q_{i_0}=q_{\mathrm{acc}}\).  The states
\(q_0,\ldots,q_{i_0-1}\) are not halting, so \(c_{i_0}\) is the halting
configuration of the computation, and the computation reaches
\(q_{\mathrm{acc}}\).

\emph{Step 3: an accepting run gives a finite model.}  Suppose the
computation reaches \(q_{\mathrm{acc}}\) at step \(s\), so that
\(q_0,\ldots,q_{s-1}\) are not halting and \(q_s=q_{\mathrm{acc}}\).  Put
\(N=2s+k\) and, for \(t\leq N\), \(\hat t=\min(t,s)\).  Define \(\mathcal M\) on
\(\{0,\ldots,N\}\) by letting \(<\) be the usual order of integers,
\[
 \mathsf{St}_q^{\mathcal M}=\{t:q_{\hat t}=q\},\qquad
 \mathsf{Hd}^{\mathcal M}=\{(t,h_{\hat t}+s):t\leq N\},
\]
\[
 \mathsf{Cl}_a^{\mathcal M}=\{(t,p):T_{\hat t}(p-s)=a\}.
\]
Real cell \(i\) is thus element \(i+s\).  Since
\(\lvert h_{\hat t}\rvert\leq\hat t\leq s\), the head element \(h_{\hat t}+s\)
lies in \(\{0,\ldots,2s\}\subseteq\{0,\ldots,N\}\).  We check the clauses.

The order clause holds for the integers, and by Step~1
\(\mathsf{Nx}(x,y)\) holds iff \(y=x+1\), and \(\mathsf{Fst}(x)\) holds iff
\(x=0\).  The determinacy clause holds because each moment has one state, one
head position and a tape function.  Start: \(q_{\hat 0}=q_0\), and the head at
moment \(0\) is element \(s\).  If \(k=0\), every \(T_w(p-s)\) is \(\sqcup\).
If \(k>0\), take \(y_i=s+i\) for \(i<k\); these lie in \(\{0,\ldots,N\}\) since
\(s+k-1\leq N\), consecutive ones satisfy \(\mathsf{Nx}\), and
\(T_w(i)=a_i\) gives \(\mathsf{Cl}_{a_i}(0,s+i)\).  For any other \(u\),
\(u-s\notin\{0,\ldots,k-1\}\), so \(T_w(u-s)=\sqcup\).

Steps: suppose \(\mathsf{St}_q(t)\), \(\mathsf{Hd}(t,p)\) and
\(\mathsf{Cl}_a(t,p)\) with \(q\) not halting and \(\delta(q,a)=(q',b,d)\).
Then \(q_{\hat t}=q\) is not halting, so \(\hat t<s\) and hence \(t=\hat t<s\).
Also \(p=h_t+s\) and \(a=T_t(h_t)\), so \(c_{t+1}=(q',T_{t+1},h_{t+1})\) is the
successor of \(c_t\).  Take \(t'=t+1\leq s\leq N\) and \(p'=h_{t+1}+s\), which
lies in \(\{0,\ldots,2s\}\).  Then \(\mathsf{Nx}(t,t')\); \(\mathsf{Mv}_d(p,p')\)
because \(p'=p+1\) for \(d=R\) and \(p'=p-1\) for \(d=L\);
\(\mathsf{St}_{q'}(t')\) because \(\widehat{t+1}=t+1\) and \(q_{t+1}=q'\);
\(\mathsf{Hd}(t',p')\) by definition; \(\mathsf{Cl}_b(t',p)\) because
\(T_{t+1}(p-s)=T_{t+1}(h_t)=b\); and if \(u\neq p\) and \(\mathsf{Cl}_c(t,u)\),
then \(T_t(u-s)=c\) and \(u-s\neq h_t\), so \(T_{t+1}(u-s)=c\) and
\(\mathsf{Cl}_c(t',u)\).

Acceptance: \(\mathsf{St}_{q_{\mathrm{acc}}}(s)\) holds.  Every
\(q_{\hat t}\) is one of \(q_0,\ldots,q_s\); the first \(s\) of these are not
halting and \(q_s=q_{\mathrm{acc}}\neq q_{\mathrm{rej}}\), so no element
satisfies \(\mathsf{St}_{q_{\mathrm{rej}}}\).

Thus \(\mathcal M\) is a model of \(\psi_{\mathcal T,w}\) with \(N+1=2s+k+1\)
elements.  Steps~2 and~3 together prove the equivalence.

\emph{Step 4: computability.}  For fixed \(\mathcal T\), the order,
determinacy, step and acceptance clauses form one fixed sentence, since \(Q\),
\(\mathcal A_{\mathrm{tape}}\) and \(\delta\) are finite.  The start clause is
assembled from \(w\) by a loop over its letters, writing \(k\) quantifiers,
\(k-1\) conjuncts \(\mathsf{Nx}(y_i,y_{i+1})\), \(k\) conjuncts
\(\mathsf{Cl}_{a_i}(t,y_i)\) and the final universal conjunct.  Composing this
with the fixed effective coding gives a computable map
\(w\mapsto\psi_{\mathcal T,w}\).
\end{proof}
\end{annexitem}

\begin{annexitem}{annex:p02-w02-p02-036-gadget}{thm:p02-w02-p02-036-gadget}{One binary relation preserves satisfiability and finite satisfiability}
\begin{proof}
Call a vertex of an \(\{E\}\)-structure \emph{looped} if it has an \(E\)-edge to
itself.  Since all the \(z_i\) in \(\pi_j(x,y)\) are bound and differ from
\(x\) and \(y\), the free variables of \(\pi_j(x,y)\) are among \(x,y\), and
by \cref{lem:fol-coincidence} its truth under an assignment depends only on
the values of \(x\) and \(y\).

\emph{Item 1.}  Fix \(\mathcal G\) with a looped vertex, so that
\(D(\mathcal G)\) is defined.  An assignment \(s\) whose values are looped
vertices is also an assignment in \(D(\mathcal G)\).  We prove the claim for
all such \(s\) by induction on the complexity of \(\varphi\).  For
\(x=y\) both sides say \(s(x)=s(y)\), since equality is identity in both
structures.  For \(R_j(x,y)\), \(\mathcal G,s\entails\pi_j(x,y)\) iff
\((s(x),s(y))\in R_j^{D(\mathcal G)}\) by the definition of
\(D(\mathcal G)\), that is, iff \(D(\mathcal G),s\entails R_j(x,y)\); the
case \(P_i(x)\) is the same with \(\pi_{l+i}(x,x)\).  The formula \(\bottom\)
is false on both sides.  The connective cases follow from the induction
hypotheses, since the translation commutes with the connectives and both
sides use the same truth tables.  For \(\forall x\,\varphi\):
\(\mathcal G,s\entails\forall x\,(Exx\to\varphi^{I})\) iff
\(\mathcal G,s[x\mapsto v]\entails\varphi^{I}\) for every looped vertex \(v\),
iff, by the induction hypothesis applied to the assignments
\(s[x\mapsto v]\), whose values are looped, \(D(\mathcal G),s[x\mapsto v]
\entails\varphi\) for every \(v\in\operatorname{dom}D(\mathcal G)\), iff
\(D(\mathcal G),s\entails\forall x\,\varphi\).  The case \(\exists x\,\varphi\)
is the same with ``some'' in place of ``every''.

\emph{Item 2.}  Let \(\mathcal M\) be an \(\Lang\)-structure and
\(\mathcal G=G(\mathcal M)\).  The vertices \(c^f_i\) are called
\emph{auxiliary}.  Every edge leaving an auxiliary vertex goes to a different
vertex (\(c^f_i\to c^f_{i+1}\) or \(c^f_j\to b\)), so no auxiliary vertex is
looped, and every element of \(\operatorname{dom}\mathcal M\) is looped.  Hence
\(D(\mathcal G)\) has domain \(\operatorname{dom}\mathcal M\).  It remains to
show, for elements \(a,b\) and an index \(j\), that
\(\mathcal G\entails\pi_j[a,b]\) iff \((a,b,j)\) is a fact; this gives
\(R_j^{D(\mathcal G)}=R_j^{\mathcal M}\) and
\(P_i^{D(\mathcal G)}=P_i^{\mathcal M}\).

If \(f=(a,b,j)\) is a fact, then \(z_i\mapsto c^f_i\) witnesses \(\pi_j\):
the edges \(a\to c^f_1\), \(c^f_i\to c^f_{i+1}\) and \(c^f_j\to b\) exist, and
the \(c^f_i\) are not looped.  Conversely, let \(v_1,\ldots,v_j\) be unlooped
vertices with edges \(a\to v_1\), \(v_i\to v_{i+1}\) and \(v_j\to b\).  The
vertex \(v_1\) is unlooped, so it is auxiliary.  The edges into auxiliary vertices are
\(a'\to c^{f'}_1\), with \(a'\) the first coordinate of \(f'\), and
\(c^{f'}_{i}\to c^{f'}_{i+1}\), whose source is auxiliary.  Since \(a\) is not
auxiliary, \(v_1=c^{f'}_1\) for a fact \(f'=(a,b',j')\).  Each \(c^{f'}_i\) has
exactly one outgoing edge, to \(c^{f'}_{i+1}\) if \(i<j'\) and to \(b'\) if
\(i=j'\).  Induction on \(i\) gives \(v_i=c^{f'}_i\) for
\(i\leq\min(j,j')\).  If \(j'<j\), then \(v_{j'+1}\) is the successor of
\(c^{f'}_{j'}\), namely \(b'\), which is looped; this is a contradiction.  If
\(j'>j\), then \(b\) is the successor of \(v_j=c^{f'}_j\), namely the
auxiliary vertex \(c^{f'}_{j+1}\), whereas \(b\) is an element; this is
also a contradiction.  So \(j'=j\), and \(b\) is the successor of
\(c^{f'}_j\), namely \(b'\).  Hence \((a,b,j)=f'\) is a fact.

If \(\mathcal M\) is finite, it has finitely many facts, each contributing at
most \(l+m\) auxiliary vertices, so \(G(\mathcal M)\) is finite.

\emph{Item 3.}  If \(\mathcal M\) is a finite model of \(\varphi\), then
\(G(\mathcal M)\) is finite.  Its domain elements are looped vertices, and the
domain is nonempty, so \(G(\mathcal M)\entails\exists x\,Exx\).  By item~2,
\(D(G(\mathcal M))=\mathcal M\entails\varphi\), so by item~1
\(G(\mathcal M)\entails\varphi^{I}\).  Hence
\(G(\mathcal M)\entails\varphi^{\sharp}\).  Conversely, if \(\mathcal G\) is a
finite model of \(\varphi^{\sharp}\), it has a looped vertex, and by item~1
\(D(\mathcal G)\entails\varphi\).  The structure \(D(\mathcal G)\) is finite
because its domain is a subset of \(\operatorname{dom}\mathcal G\).  The
translation is defined by recursion on formulas with explicitly written
clauses, so \(\varphi\mapsto\varphi^{\sharp}\) is computable.

\emph{Item 4.}  The argument of item~3 without the finiteness clauses applies.
If \(\mathcal M\) is any model of \(\varphi\), then \(G(\mathcal M)\) has a
looped vertex and \(G(\mathcal M)\entails\varphi^{\sharp}\) by items 1 and~2.
If \(\mathcal G\) is any model of \(\varphi^{\sharp}\), then \(\mathcal G\) has
a looped vertex and \(D(\mathcal G)\entails\varphi\) by item~1.  So \(\varphi\)
is satisfiable iff \(\varphi^{\sharp}\) is.  Together with item~3, \(\varphi\)
and \(\varphi^{\sharp}\) agree on satisfiability and on finite
satisfiability.  Each of the three classes is determined by these two
properties: \(\mathrm{Fin\text{-}sat}\) by finite satisfiability,
\(\mathrm{Non\text{-}sat}\) by unsatisfiability, and
\(\mathrm{Inf\text{-}axioms}\) by satisfiability without finite
satisfiability.  So \(\varphi\) and \(\varphi^{\sharp}\) lie in corresponding
classes.
\end{proof}
\end{annexitem}

\begin{annexitem}{annex:p02-w02-p02-036-transfer}{thm:p02-w02-p02-036-transfer}{Transfer to any signature with a binary symbol}
\begin{proof}
For an \(\{E\}\)-structure \(\mathcal G\), let \(\mathcal G^R\) be the
\(\{R\}\)-structure with the same domain and \(R^{\mathcal G^R}=E^{\mathcal G}\).
The satisfaction clauses for \(\chi\) in \(\mathcal G\) and for \(\chi[R/E]\)
in \(\mathcal G^R\) are the same clauses applied to the same relation, so
induction on complexity gives \(\mathcal G,s\entails\theta\) iff
\(\mathcal G^R,s\entails\theta[R/E]\) for every \(\{E\}\)-formula \(\theta\) and
assignment \(s\).

Suppose \(\mathcal G\) is a finite model of \(\chi\).  Let \(\mathcal M'\) be
the \(\Lang'\)-structure with domain \(\operatorname{dom}\mathcal G\),
\(R^{\mathcal M'}=E^{\mathcal G}\), and every other relation symbol
interpreted by the empty relation of its arity.  Its reduct to \(\{R\}\) is
\(\mathcal G^R\), which satisfies \(\chi[R/E]\).  By
\cref{lem:reduct-invariance}, \(\mathcal M'\entails\chi[R/E]\), and
\(\mathcal M'\) is finite.  Conversely, if \(\mathcal M'\) is a finite model of
\(\chi[R/E]\), its reduct to \(\{R\}\) satisfies \(\chi[R/E]\) by
\cref{lem:reduct-invariance}.  That reduct is \(\mathcal G^R\) for the
\(\{E\}\)-structure \(\mathcal G\) with the same domain and
\(E^{\mathcal G}=R^{\mathcal M'}\), so \(\mathcal G\) is a finite model of
\(\chi\).

The same two arguments, with ``finite'' deleted throughout, show that
\(\chi\) has a model iff \(\chi[R/E]\) has one: the expansion
\(\mathcal M'\) of a model \(\mathcal G^R\) is a model of \(\chi[R/E]\), and
the reduct of a model of \(\chi[R/E]\) to \(\{R\}\), read back as an
\(\{E\}\)-structure, is a model of \(\chi\).  The domain is the same in every
step, so finiteness is preserved in both directions as well.
\end{proof}
\end{annexitem}

\begin{annexitem}{annex:p02-w02-p02-036-trakhtenbrot}{thm:p02-w02-p02-036-trakhtenbrot}{Proof of Trakhtenbrot's theorem}
\begin{proof}
Let \(R\) be a binary relation symbol of \(\Lang\).

\emph{Enumerability.}  \(\mathrm{Fin\text{-}sat}(X)\) is computably enumerable,
and the set of finitely valid sentences is co-computably enumerable, by items
2 and~4 of Lemma~\ref{thm:p02-w02-p02-036-finsat-ce}.

\emph{Undecidability.}  The set \(K\) is computably enumerable
(\cref{thm:halting-undecidable}), so some Turing machine \(\mathcal R\)
recognizes it (\cref{thm:ce-equivalent-forms}).  Its input alphabet contains
the digits \(0,1\), since it is run on binary numerals, and its computation on
the input word \(\operatorname{bin}(e)\) reaches \(q_{\mathrm{acc}}\) iff
\(e\in K\).  The
run signature \(\Lang_{\mathcal R}\) consists of binary and unary relation
symbols, so Definition~\ref{def:p02-w02-p02-036-gadget} applies to it.  Put
\[
 g(e)\defeq\bigl(\psi_{\mathcal R,\operatorname{bin}(e)}\bigr)^{\sharp}[R/E].
\]
The map \(g\) is computable: \(e\mapsto\operatorname{bin}(e)\) is computable,
and so are the three syntactic maps by
Lemmas~\ref{thm:p02-w02-p02-036-run-sentence}
and~\ref{thm:p02-w02-p02-036-gadget} and by the letter-for-letter renaming.
By Lemmas~\ref{thm:p02-w02-p02-036-run-sentence},
\ref{thm:p02-w02-p02-036-gadget}(3) and~\ref{thm:p02-w02-p02-036-transfer},
\[
\begin{aligned}
 e\in K
 &\iff \psi_{\mathcal R,\operatorname{bin}(e)}\in
      \mathrm{Fin\text{-}sat}(\operatorname{Sent}(\Lang_{\mathcal R}))\\
 &\iff \psi_{\mathcal R,\operatorname{bin}(e)}^{\sharp}\in
      \mathrm{Fin\text{-}sat}(\operatorname{Sent}(\{E\}))\\
 &\iff g(e)\in\mathrm{Fin\text{-}sat}(X).
\end{aligned}
\]
If \(\mathrm{Fin\text{-}sat}(X)\) were decidable, computing \(g(e)\) and then
deciding membership would decide \(K\), against
\cref{thm:halting-undecidable}.

\emph{Finite validity is not computably enumerable.}  Suppose the set \(F\) of
finitely valid sentences in \(X\) were computably enumerable.  For
\(\varphi\in X\), \(\varphi\notin\mathrm{Fin\text{-}sat}(X)\) iff no finite
structure satisfies \(\varphi\) iff \(\neg\varphi\in F\).  A machine that
checks \(\varphi\in X\) and runs a recognizer for \(F\) on \(\neg\varphi\)
then recognizes \(X\setminus\mathrm{Fin\text{-}sat}(X)\).  The complement of
\(\mathrm{Fin\text{-}sat}(X)\) in \(\mathbb N\) is the union of this set with
the decidable set \(\mathbb N\setminus X\), so it is computably enumerable.
Since \(\mathrm{Fin\text{-}sat}(X)\) is computably enumerable as well, the
complement criterion makes it decidable, contradicting the previous
paragraph.
\end{proof}
\end{annexitem}

\begin{annexitem}{annex:p02-w02-p02-036-classes}{thm:p02-w02-p02-036-classes}{Effective standing of the countermodel classes}
\begin{proof}
Let \(R\) be a binary relation symbol of \(\Lang\).

\emph{Item 1.}  \(\mathrm{Fin\text{-}sat}(X)\) and
\(\mathrm{Non\text{-}sat}(X)\) are computably enumerable by items 2 and~3 of
Lemma~\ref{thm:p02-w02-p02-036-finsat-ce}, so their union is.  Since the three
classes partition \(X\), the union is \(X\setminus\mathrm{Inf\text{-}axioms}(X)\).

\emph{Item 2.}  Nonempty: let \(\chi\) be the result of replacing \(K\) by
\(R\) in \(\prec_K\).  By Lemma~\ref{thm:p02-w02-p02-036-infinity-axiom},
\(\prec_K\) has a model and no finite model.  Lemma~\ref{thm:p02-w02-p02-036-transfer},
applied with the binary symbol \(K\) in the role of \(E\), gives that \(\chi\)
has a model and no finite model.  So \(\chi\in\mathrm{Inf\text{-}axioms}(X)\).

Not decidable: suppose a machine decides \(\mathrm{Inf\text{-}axioms}(X)\).
Decide \(\mathrm{Fin\text{-}sat}(X)\) as follows.  Given \(\varphi\), reject if
\(\varphi\notin X\) or \(\varphi\in\mathrm{Inf\text{-}axioms}(X)\).  Otherwise
\(\varphi\) lies in exactly one of \(\mathrm{Fin\text{-}sat}(X)\) and
\(\mathrm{Non\text{-}sat}(X)\).  Run recognizers for both alternately, one step
at a time.  Exactly one of them eventually accepts.  Accept if the first one
does and reject if the second one does.  This procedure halts on every input
and decides \(\mathrm{Fin\text{-}sat}(X)\), contradicting
Corollary~\ref{thm:p02-w02-p02-036-trakhtenbrot}.

Not computably enumerable: the complement of \(\mathrm{Inf\text{-}axioms}(X)\)
in \(\mathbb N\) is the union of \(\mathbb N\setminus X\) and the set of item~1,
so it is computably enumerable.  If \(\mathrm{Inf\text{-}axioms}(X)\) were
computably enumerable too, the complement criterion would make it decidable,
which was just excluded.

\emph{Item 3.}  For \(\varphi\in X\): \(\varphi\) is finitely valid iff
\(\neg\varphi\) has no finite model, and \(\varphi\) is not valid iff
\(\neg\varphi\) has a model.  So \(\varphi\) is finitely valid but not valid
iff \(\neg\varphi\in\mathrm{Inf\text{-}axioms}(X)\).  In that case the
countermodels of \(\varphi\), that is, the models of \(\neg\varphi\), exist and
are all infinite.  The sentence \(\neg\chi\), with \(\chi\) from item~2,
belongs to the set.  Suppose the set \(F'\) of finitely valid but not valid
sentences in \(X\) were computably enumerable.  The sentences \(\psi\) and
\(\neg\neg\psi\) have the same models, so for \(\psi\in X\),
\(\psi\in\mathrm{Inf\text{-}axioms}(X)\) iff
\(\neg\neg\psi\in\mathrm{Inf\text{-}axioms}(X)\) iff \(\neg\psi\in F'\).  Then
checking \(\psi\in X\) and running a recognizer for \(F'\) on \(\neg\psi\) would
recognize \(\mathrm{Inf\text{-}axioms}(X)\), contradicting item~2.
\end{proof}
\end{annexitem}

\begin{annexitem}{annex:p02-w02-p02-036-fmp}{thm:p02-w02-p02-036-fmp}{The finite model property gives decidability}
\begin{proof}
Let the symbols of \(\Lang\) be relation symbols \(R_i\) of arity \(m_i\),
function symbols \(f_j\) of arity \(n_j\geq1\) and constant symbols \(c_k\),
finitely many of each kind.

\emph{Step 1: finite model checking.}  Let \(\mathfrak A\) have domain
\(\{0,\ldots,n-1\}\), given as in
Convention~\ref{def:p02-w02-p02-036-finite-data} by the lists of the
relations, the tables of the functions and the values of the constants.  For
a term \(t\) and a finite map \(s_0\) from a finite set of variables
containing the variables of \(t\) into \(\{0,\ldots,n-1\}\), compute
\(V(t,s_0)\) by recursion on \(t\): a variable \(x\) gets \(s_0(x)\), a
constant \(c_k\) gets its given value, and \(f_j(t_1,\ldots,t_{n_j})\) gets the
table entry of \(f_j\) at \((V(t_1,s_0),\ldots,V(t_{n_j},s_0))\).  These are the
clauses of the term value in \cref{sec:fol-semantics-01}, so by induction on
\(t\), \(V(t,s_0)\) is the value of \(t\) under every assignment extending
\(s_0\).  For formulas, compute a truth value \(V(\theta,s_0)\) exactly as in
item~1 of the proof of Lemma~\ref{thm:p02-w02-p02-036-finsat-ce}, with two
changes in the atomic case: \(t_1=t_2\) is true iff \(V(t_1,s_0)=V(t_2,s_0)\),
and \(R_i(t_1,\ldots,t_{m_i})\) is true iff the tuple of the values
\(V(t_l,s_0)\) occurs in the list for \(R_i\).  The recursion terminates and
agrees with satisfaction, by the argument given there together with the term
clause just proved.  So truth of a sentence in a finite structure given in
this format is decidable.

\emph{Step 2: finite satisfiability of sentences is computably enumerable.}
For each \(n\geq1\) there are finitely many \(\Lang\)-structures with domain
\(\{0,\ldots,n-1\}\): finitely many choices of each relation, of each table,
and of each constant, and they can be listed effectively.  A machine that, on
a sentence \(\chi\in\operatorname{Sent}(\Lang)\), runs through
\(n=1,2,\ldots\) and through these structures and accepts when Step~1 reports a
model, accepts exactly the sentences with a model on some
\(\{0,\ldots,n-1\}\).  Every finite \(\Lang\)-structure is isomorphic to one
with such a domain: transport the interpretations along a bijection from its
domain onto \(\{0,\ldots,n-1\}\).  By \cref{lem:iso-invariance}, the machine
therefore recognizes \(\mathrm{Fin\text{-}sat}(\operatorname{Sent}(\Lang))\).

\emph{Step 3: from formulas to sentences.}  For a formula \(\varphi\) whose free
variables, listed in their order of first occurrence, are \(x_1,\ldots,x_p\),
let \(\varphi^{\exists}\defeq\exists x_1\cdots\exists x_p\,\varphi\).  The free
variables are computed by recursion on \(\varphi\), so
\(\varphi\mapsto\varphi^{\exists}\) is computable.  By the satisfaction clause
for \(\exists\) and \cref{lem:fol-coincidence}, a structure \(\mathfrak A\)
satisfies \(\varphi^{\exists}\) iff \(\mathfrak A,s\entails\varphi\) for some
assignment \(s\).  Hence \(\varphi\) is satisfiable iff \(\varphi^{\exists}\)
is, and \(\varphi\) has a finite model iff \(\varphi^{\exists}\) has one.

\emph{Step 4: unsatisfiability is computably enumerable.}  A sentence \(\chi\)
has no model iff \(\neg\chi\) is true in every structure for a signature
containing the symbols of \(\chi\).  By \cref{lem:reduct-invariance}, as in
item~3 of the proof of Lemma~\ref{thm:p02-w02-p02-036-finsat-ce}, this is the
validity of \(\neg\chi\), and the valid sentences form a computably enumerable
set (\cref{thm:fol-validity-enumerable}).

\emph{Step 5: decision.}  By the finite model property,
\(\mathrm{Sat}(X)=\mathrm{Fin\text{-}sat}(X)\).  By Steps 2 and~3, a machine
that checks \(\varphi\in X\), computes \(\varphi^{\exists}\) and runs the
recognizer of Step~2 on it recognizes \(\mathrm{Sat}(X)\).  By Steps 3 and~4,
a machine that checks \(\varphi\in X\) and runs a recognizer for the valid
sentences on \(\neg\varphi^{\exists}\) recognizes \(X\setminus\mathrm{Sat}(X)\).
The complement of \(\mathrm{Sat}(X)\) in \(\mathbb N\) is the union of
\(X\setminus\mathrm{Sat}(X)\) and the decidable set \(\mathbb N\setminus X\), so
it is computably enumerable.  Both \(\mathrm{Sat}(X)\) and its complement are
computably enumerable, so \(\mathrm{Sat}(X)\) is decidable by the complement
criterion.
\end{proof}
\end{annexitem}

\begin{annexitem}{annex:p02-w02-p02-036-halting}{thm:p02-w02-p02-036-halting}{The three halting problems are pairwise recursively inseparable}
\begin{proof}
By the symmetry remark after
Definition~\ref{def:p02-w02-p02-036-inseparable}, it suffices to show that
no decidable set separates \(H_2\) from \(H_1\), \(H_2\) from \(H_\infty\), or
\(H_1\) from \(H_\infty\).  Three facts about machines and codes are used.

\emph{(a) Coding a machine.}  Let \(N\) be a Turing machine whose input
alphabet contains the digits \(0,1\).  The normalization in the proof of
\cref{thm:enumeration}(b), a surrogate start state if the start state is
halting followed by a renaming onto the designated numbering, produces a valid
code \(n\) such that, by \cref{lem:renaming-invariance}, \(M_n\) and \(N\)
accept together, reject together and diverge together on every word over the
digits.  The construction operates on the finite transition table, so when
\(N\) is obtained from a decoded machine by a fixed finite modification, the
resulting code is a computable function of the original code.

\emph{(b) Writing the code first.}  For \(e\in V\) let
\(\operatorname{bin}(e)=b_0b_1\cdots b_{L-1}\), with \(L\geq1\).  Let \(N_e\) have
the states of \(M_e\) together with fresh states \(p_0,\ldots,p_{L-1}\) and
\(r_1,\ldots,r_L\), start state \(p_0\), the halting states of \(M_e\), the
alphabets of \(M_e\), the transitions of \(M_e\), and, for every tape symbol
\(a\),
\[
\begin{aligned}
 \delta(p_i,a)&=(p_{i+1},b_i,R)\quad(i<L-1),&\qquad
 \delta(p_{L-1},a)&=(r_L,b_{L-1},R),\\
 \delta(r_j,a)&=(r_{j-1},a,L)\quad(2\leq j\leq L),&\qquad
 \delta(r_1,a)&=(q_0^{e},a,L),
\end{aligned}
\]
where \(q_0^{e}\) is the start state of \(M_e\).  On the empty input, \(N_e\)
writes \(b_i\) in cell \(i\) during step \(i+1\) for \(i<L\), arriving in state
\(r_L\) at cell \(L\); then it moves left \(L\) times without writing, passing
through \(r_L,\ldots,r_1\), and after \(2L\) steps it is in state \(q_0^{e}\) at
cell \(0\) with tape \(T_{\operatorname{bin}(e)}\).  That is the initial
configuration of \(M_e\) on \(\operatorname{bin}(e)\), and none of the \(2L\)
earlier states is halting.  From then on \(N_e\) uses only transitions of
\(M_e\), so the two computations agree.  Let \(w(e)\) be the code given by (a)
for \(N_e\); \(w\) is computable on \(V\).  Then \(M_{w(e)}\) on the empty input
accepts, rejects or diverges exactly as \(M_e\) does on \(\operatorname{bin}(e)\).

\emph{(c) Turning one halting state into a loop.}  For \(e\in V\) let
\(N^{\mathrm{rej}}_e\) be \(M_e\) with one fresh state \(q^{*}\) made the
rejecting state, and with the old rejecting state \(q_{\mathrm{rej}}\) made an
ordinary state with \(\delta(q_{\mathrm{rej}},a)=(q_{\mathrm{rej}},a,R)\) for
every \(a\).  No transition leads to \(q^{*}\).  On the empty input the two
computations agree until \(M_e\) reaches \(q_{\mathrm{rej}}\); from there
\(N^{\mathrm{rej}}_e\) moves right forever.  So \(N^{\mathrm{rej}}_e\) accepts
iff \(M_e\) accepts, and diverges iff \(M_e\) rejects or diverges.  Let
\(\lambda_{\mathrm{rej}}(e)\) be its code by (a).  Define
\(N^{\mathrm{acc}}_e\) and \(\lambda_{\mathrm{acc}}(e)\) in the same way with
the roles of the accepting and rejecting states exchanged: it rejects iff
\(M_e\) rejects, and diverges iff \(M_e\) accepts or diverges.  Both maps are
computable on \(V\).

\emph{\(H_2\) from \(H_1\).}  Let
\(D_1\defeq\{e\in V:M_e\text{ rejects }\operatorname{bin}(e)\}\) and
\(D_2\defeq\{e\in V:M_e\text{ accepts }\operatorname{bin}(e)\}\).  Suppose
first that a decidable set \(W\) separates \(D_2\) from \(D_1\).  By the machine
form of decidability (\cref{sec:turing-machines-03}) some machine decides
\(W\); it halts on every binary numeral, so its input alphabet contains both
digits.  Exchange its accepting and rejecting states to obtain \(N\), which
on \(\operatorname{bin}(e)\) rejects if \(e\in W\) and accepts if \(e\notin W\).
Let \(n\) be a code for \(N\) by (a).  If \(n\in W\), then \(M_n\) rejects
\(\operatorname{bin}(n)\), so \(n\in D_1\), against \(W\cap D_1=\varnothing\).
If \(n\notin W\), then \(M_n\) accepts \(\operatorname{bin}(n)\), so
\(n\in D_2\subseteq W\), a contradiction.  So no decidable set separates
\(D_2\) from \(D_1\).  Now suppose a decidable \(W\) separates \(H_2\) from
\(H_1\), and put \(W'\defeq\{e\in V:w(e)\in W\}\), which is decidable since
\(V\) is decidable and \(w\) is computable.  By (b), \(e\in D_2\) gives
\(w(e)\in H_2\subseteq W\), and \(e\in D_1\) gives \(w(e)\in H_1\), hence
\(w(e)\notin W\).  So \(W'\) separates \(D_2\) from \(D_1\), which was excluded.

\emph{\(H_2\) from \(H_\infty\).}  Suppose a decidable \(W\) separates \(H_2\)
from \(H_\infty\), and put \(W'\defeq\{e\in V:\lambda_{\mathrm{rej}}(e)\in W\}\).
By (c), \(e\in H_2\) gives \(\lambda_{\mathrm{rej}}(e)\in H_2\subseteq W\), and
\(e\in H_1\) gives \(\lambda_{\mathrm{rej}}(e)\in H_\infty\), hence
\(\lambda_{\mathrm{rej}}(e)\notin W\).  So the decidable set \(W'\) separates
\(H_2\) from \(H_1\), which was excluded.

\emph{\(H_1\) from \(H_\infty\).}  Suppose a decidable \(W\) separates \(H_1\)
from \(H_\infty\), and put \(W'\defeq\{e\in V:\lambda_{\mathrm{acc}}(e)\in W\}\).
By (c), \(e\in H_1\) gives \(\lambda_{\mathrm{acc}}(e)\in H_1\subseteq W\), and
\(e\in H_2\) gives \(\lambda_{\mathrm{acc}}(e)\in H_\infty\), hence
\(\lambda_{\mathrm{acc}}(e)\notin W\).  So \(W'\) separates \(H_1\) from \(H_2\),
and its complement separates \(H_2\) from \(H_1\), which was excluded.
\end{proof}
\end{annexitem}

\begin{annexitem}{annex:p02-w02-p02-036-simulation}{thm:p02-w02-p02-036-simulation}{Simulation lemma}
\begin{proof}
Write \(K\), \(S\), and so on for the interpretations in \(\mathfrak A\).

\emph{Order facts.}  By \(\mathrm{ORDER}^{-}\), \(K\) is a strict linear order
and \(S(x,y)\) holds iff \(y\) is an immediate \(K\)-successor of \(x\).  In
particular \(S(x,y)\) implies \(K(x,y)\).  An element has at most one
\(S\)-successor: if \(S(x,y)\), \(S(x,y')\) and \(y\neq y'\), say \(K(y,y')\) by
linearity, then \(y\) lies strictly between \(x\) and \(y'\), against the
immediacy of \(y'\).  Likewise an element has at most one
\(S\)-predecessor.

\emph{Cells.}  Put \(\sigma_0\defeq z\).  For \(x\geq0\), if \(\sigma_x\) is
defined and has an \(S\)-successor, let \(\sigma_{x+1}\) be that successor; for
\(x\leq0\), if \(\sigma_x\) is defined and has an \(S\)-predecessor, let
\(\sigma_{x-1}\) be that predecessor.  The set \(I\) of indices where
\(\sigma\) is defined is an interval of \(\mathbb Z\) containing \(0\).  If
\(x,x+1\in I\), then \(S(\sigma_x,\sigma_{x+1})\): for \(x\geq0\) by the
definition of \(\sigma_{x+1}\), and for \(x+1\leq0\) by the definition of
\(\sigma_x\).  Hence \(K(\sigma_x,\sigma_{x+1})\), and by transitivity
\(K(\sigma_x,\sigma_y)\) whenever \(x<y\) in \(I\).  By irreflexivity \(\sigma\)
is injective.

\emph{Times.}  We define \(\tau_t\) and prove items 1--3 by induction on \(t\),
for every \(t\) such that \(c_t\) exists, together with
\(S(\tau_{t-1},\tau_t)\) for \(t\geq1\).

For \(t=0\), put \(\tau_0\defeq z\).  The computation starts in
\(c_0=(q_0,T_{\varepsilon},0)\), where \(T_{\varepsilon}\) is blank everywhere.
\(\mathrm{START}\) at \(x\mapsto z\) gives \(I_{q_0}(z)\), \(H(z,z)\) and
\(T_{\sqcup}(z,y)\) for every element \(y\).  This is item~1, item~2 with
\(h_0=0\in I\) and \(\sigma_0=z\), and item~3, since \(B_0(x)=\sqcup\) for every
\(x\).

Suppose the claims hold for \(t\) and \(c_{t+1}\) exists.  Then \(q_t\) is not
halting, so by item~1 \(\mathrm{Act}(\tau_t)\) holds.  Apply \(\prec_{K,M}\) at
\(x\mapsto\tau_t\), with any element for \(y\): it gives an \(S\)-successor
of \(\tau_t\), and we let \(\tau_{t+1}\) be this unique successor.  Apply
\(\prec_{K,M}\) at \(x\mapsto\sigma_{h_t}\) with \(y\mapsto\tau_t\): since
\(H(\tau_t,\sigma_{h_t})\) holds by item~2, \(\sigma_{h_t}\) has an
\(S\)-successor and an \(S\)-predecessor.  Hence \(h_t+1\in I\): if
\(h_t\geq0\), \(\sigma_{h_t+1}\) is that successor, and if \(h_t<0\), then
\(h_t+1\) lies between \(h_t\) and \(0\).  Symmetrically \(h_t-1\in I\).

Let \(a=B_t(h_t)\) and \(\delta(q_t,a)=(q',b,d)\), so that
\(c_{t+1}=(q',B_{t+1},h_{t+1})\) with \(B_{t+1}(h_t)=b\), \(B_{t+1}(x)=B_t(x)\)
for \(x\neq h_t\), and \(h_{t+1}=h_t\pm1\).  By items 1--3 for \(t\),
\(I_{q_t}(\tau_t)\), \(H(\tau_t,\sigma_{h_t})\) and \(T_a(\tau_t,\sigma_{h_t})\)
hold.  If \(d=R\), the right-move clause of \(\mathrm{STEP}_M\) for \((q_t,a)\)
with \(t\mapsto\tau_t\), \(x\mapsto\sigma_{h_t}\), \(t'\mapsto\tau_{t+1}\),
\(x'\mapsto\sigma_{h_t+1}\) has a true antecedent, since
\(S(\tau_t,\tau_{t+1})\) and \(S(\sigma_{h_t},\sigma_{h_t+1})\).  It gives
\(I_{q'}(\tau_{t+1})\), \(H(\tau_{t+1},\sigma_{h_t+1})\) and
\(T_b(\tau_{t+1},\sigma_{h_t})\).  If \(d=L\), the left-move clause with
\(x\mapsto\sigma_{h_t-1}\) and \(x'\mapsto\sigma_{h_t}\) has a true antecedent,
since \(S(\sigma_{h_t-1},\sigma_{h_t})\), and gives \(I_{q'}(\tau_{t+1})\),
\(H(\tau_{t+1},\sigma_{h_t-1})\) and \(T_b(\tau_{t+1},\sigma_{h_t})\).  In
both cases items 1 and~2 hold for \(t+1\), and \(T_{B_{t+1}(h_t)}\) holds at
\((\tau_{t+1},\sigma_{h_t})\).

For any element \(y\neq\sigma_{h_t}\) and any letter \(c\) with
\(T_c(\tau_t,y)\), the no-change clause with \(t\mapsto\tau_t\),
\(x\mapsto\sigma_{h_t}\), \(t'\mapsto\tau_{t+1}\) gives \(T_c(\tau_{t+1},y)\).
If \(y=\sigma_x\) with \(x\in I\) and \(x\neq h_t\), then \(y\neq\sigma_{h_t}\)
by injectivity, and with \(c=B_t(x)=B_{t+1}(x)\) this gives item~3 at
\(\sigma_x\).  If \(y\) is outside the image of \(\sigma\), then
\(y\neq\sigma_{h_t}\), and with \(c=\sqcup\) this gives item~3 at \(y\).  This
completes the induction.

Finally \(S(\tau_t,\tau_{t+1})\) gives \(K(\tau_t,\tau_{t+1})\), so
\(K(\tau_s,\tau_t)\) for \(s<t\) by transitivity, and the \(\tau_t\) are
pairwise distinct by irreflexivity.
\end{proof}
\end{annexitem}

\begin{annexitem}{annex:p02-w02-p02-036-refined}{thm:p02-w02-p02-036-refined}{Refined reduction property}
\begin{proof}
Let \(c_t=(q_t,B_t,h_t)\) be the computation of \(M\) on the empty input.
Each step moves the head by one cell, so \(\lvert h_t\rvert\leq t\).

\emph{Intended structures.}  Suppose the computation does not reach
\(q_{\mathrm{rej}}\).  If it halts, it does so in \(q_{\mathrm{acc}}\) at some
step \(k\); put \(\hat t\defeq\min(t,k)\) for \(t\geq0\).  If it does not halt,
put \(\hat t\defeq t\).  Let \(D\) be either \(\mathbb Z\), or, when the
computation halts at step \(k\), the set \(D_k\defeq\{-k-1,\ldots,k+1\}\).
Define \(\mathfrak A_D\) on the domain \(D\) by
\[
\begin{gathered}
 K=\{(x,y)\in D^2:x<y\},\qquad S=\{(x,x+1):x,x+1\in D\},\qquad Z=\{0\},\\
 I_q=\{t\in D:t\geq0,\ q_{\hat t}=q\},\qquad
 H=\{(t,h_{\hat t}):t\in D,\ t\geq0\},\\
 T_a=\{(t,x)\in D^2:t\geq0,\ B_{\hat t}(x)=a\}.
\end{gathered}
\]
For \(D=D_k\) every \(h_{\hat t}\) lies in \(D\), since
\(\lvert h_{\hat t}\rvert\leq\hat t\leq k\).  We check that
\(\mathfrak A_D\entails\rho_{M,\prec}\).
\begin{itemize}
\item \(\mathrm{ORDER}^{-}\): \(K\) is the usual strict order of an interval of
  integers, and \(y\) is an immediate successor of \(x\) there iff
  \(y=x+1\).
\item \(\mathrm{START}\): the only element of \(Z\) is \(0\).
  \(\hat0=0\), \(q_0\) is the start state, \(h_0=0\), and \(B_0\) is blank
  everywhere, so \(I_{q_0}(0)\), \(H(0,0)\) and \(T_{\sqcup}(0,y)\) for all
  \(y\in D\).
\item Right-move clause for \((q,a)\) with \(\delta(q,a)=(q',b,R)\): suppose
  \(I_q(t)\), \(H(t,x)\), \(T_a(t,x)\), \(S(t,t')\) and \(S(x,x')\).  Then
  \(t\geq0\) and \(q_{\hat t}=q\) is not halting, so the computation has not
  halted at \(\hat t\); hence \(\hat t=t\) and \(c_{t+1}\) exists.  Also
  \(x=h_t\), \(a=B_t(h_t)\), \(t'=t+1\) and \(x'=h_t+1\).  So
  \(c_{t+1}=(q',B_{t+1},h_t+1)\) with \(B_{t+1}(h_t)=b\), and
  \(\widehat{t+1}=t+1\).  This gives \(I_{q'}(t')\), \(H(t',x')\) and
  \(T_b(t',x)\).
\item Left-move clause for \((q,a)\) with \(\delta(q,a)=(q',b,L)\): suppose
  \(I_q(t)\), \(H(t,x')\), \(T_a(t,x')\), \(S(t,t')\) and \(S(x,x')\).  As
  before \(\hat t=t\), \(c_{t+1}\) exists, \(x'=h_t\), \(a=B_t(h_t)\),
  \(t'=t+1\), and now \(x=h_t-1=h_{t+1}\) and \(B_{t+1}(h_t)=b\).  This gives
  \(I_{q'}(t')\), \(H(t',x)\) and \(T_b(t',x')\).
\item No-change clause: suppose \(H(t,x)\), \(y\neq x\), \(T_a(t,y)\) and
  \(S(t,t')\).  Then \(t\geq0\), \(x=h_{\hat t}\), \(B_{\hat t}(y)=a\) and
  \(t'=t+1\).  If \(\hat t=t\) and \(c_{t+1}\) exists, then
  \(\widehat{t+1}=t+1\) and \(B_{t+1}(y)=B_t(y)=a\) because \(y\neq h_t\).
  Otherwise the computation halted at \(k\leq t\), and
  \(\widehat{t+1}=k=\hat t\).  In both cases \(T_a(t',y)\).
\item \(\mathrm{NONSTOP}\): no \(q_{\hat t}\) is \(q_{\mathrm{rej}}\), because
  the computation does not reach \(q_{\mathrm{rej}}\).
\item \(\prec_{K,M}\): let \(x\in D\).  Call \(x\) \emph{visited} if \(H(y,x)\)
  for some \(y\).  A visited \(x\) equals some \(h_{\hat y}\), so
  \(\lvert x\rvert\leq k\) in the case \(D=D_k\), and \(x+1\) and \(x-1\) lie
  in \(D\).  If \(\mathrm{Act}(x)\), then \(x\geq0\) and \(q_{\hat x}\) is not
  halting, so \(\hat x=x\) and, in the case \(D=D_k\), \(x<k\); hence
  \(x+1\in D\).  Now choose \(u\defeq x+1\) if \(x\) is visited or active, and
  \(u\defeq x\) otherwise; choose \(v\defeq x-1\) if \(x\) is visited, and
  \(v\defeq x\) otherwise.  These lie in \(D\).  For every \(y\): if
  \(H(y,x)\lor\mathrm{Act}(x)\), then \(x\) is visited or active and
  \(S(x,x+1)\) holds; if \(H(y,x)\), then \(x\) is visited and \(S(x-1,x)\)
  holds.  So both implications hold.
\end{itemize}
For \(D=\mathbb Z\) every element has an \(S\)-successor, so
\(\mathfrak A_{\mathbb Z}\entails\rho^{\infty}_M\) as well.  The structure
\(\mathfrak A_{D_k}\) has \(2k+3\) elements.

\emph{Item 1.}  If the computation reaches \(q_{\mathrm{rej}}\) at step \(k\)
and \(\mathfrak A\) were a model of \(\rho_{M,\prec}\), then \(\mathrm{START}\)
supplies \(z\in Z^{\mathfrak A}\), and
Lemma~\ref{thm:p02-w02-p02-036-simulation} gives \(I_{q_{\mathrm{rej}}}(\tau_k)\),
against \(\mathrm{NONSTOP}\).  So \(\rho_{M,\prec}\) has no model.  If the
computation does not reach \(q_{\mathrm{rej}}\), then
\(\mathfrak A_{\mathbb Z}\) is a model.

\emph{Item 2.}  If the computation accepts at step \(k\), then
\(\mathfrak A_{D_k}\) is a model with \(2k+3\) elements.  Conversely, let
\(\mathfrak A\) be a finite model.  By item~1 the computation does not reach
\(q_{\mathrm{rej}}\).  If it never halted, \(c_t\) would exist for every \(t\),
and Lemma~\ref{thm:p02-w02-p02-036-simulation} would give infinitely many
pairwise distinct elements \(\tau_t\) of \(\mathfrak A\).  So it halts, and
since it does not reject, it accepts.

\emph{Item 3.}  The three outcomes, rejecting, accepting and never halting,
are exclusive and exhaustive, and so are the three classes
\(\mathrm{Non\text{-}sat}\), \(\mathrm{Fin\text{-}sat}\) and
\(\mathrm{Inf\text{-}axioms}\).  Items 1 and~2 match the first two outcomes
with the first two classes, so the third outcome matches the third class.

\emph{Item 4.}  \(\rho^{\infty}_M\) implies \(\rho_{M,\prec}\), so it has no model
when the computation rejects, by item~1.  Otherwise
\(\mathfrak A_{\mathbb Z}\) is a model of it.  Every model \(\mathfrak A\) of
\(\rho^{\infty}_M\) is infinite: starting from any element \(a_0\), the clause
\(\forall x\,\exists y\,S(x,y)\) gives \(a_1,a_2,\ldots\) with
\(S(a_i,a_{i+1})\), hence \(K(a_i,a_{i+1})\) by \(\mathrm{ORDER}^{-}\), hence
\(K(a_i,a_j)\) for \(i<j\) by transitivity, and the \(a_i\) are pairwise
distinct by irreflexivity.  So \(\rho^{\infty}_M\) is then an infinity axiom.

\emph{Computability.}  For a given transition table, \(\rho_{M,\prec}\) is a
finite conjunction whose conjuncts are listed by loops over \(Q\),
\(\mathcal A_{\mathrm{tape}}\) and the table, and \(\rho^{\infty}_M\) adds one
fixed conjunct.  Composing with the fixed effective coding gives computable
maps.
\end{proof}
\end{annexitem}

\begin{annexitem}{annex:p02-w02-p02-036-inseparable}{thm:p02-w02-p02-036-inseparable}{Trakhtenbrot's theorem, inseparable form}
\begin{proof}
Let \(R\) be a binary relation symbol of \(\Lang\).

\emph{The reduction.}  For \(e\in V\), the signature \(\Lang_{M_e}^{\prec}\)
consists of binary and unary relation symbols, listed computably from the
decoded table: \(K\), \(S\), \(H\) and the \(T_a\) in the order of the symbol
numbers as \(R_1,\ldots,R_l\), and \(Z\) and the \(I_q\) in the order of the
state numbers as \(P_1,\ldots,P_m\).  So
Definition~\ref{def:p02-w02-p02-036-gadget} applies, and we put
\[
 g(e)\defeq\bigl(\rho_{M_e,\prec}\bigr)^{\sharp}[R/E]\in X.
\]
The map \(g\) is computable on \(V\), by
Theorem~\ref{thm:p02-w02-p02-036-refined},
Lemma~\ref{thm:p02-w02-p02-036-gadget}(3) and the letter-for-letter renaming.
By Theorem~\ref{thm:p02-w02-p02-036-refined}(1)--(3), \(\rho_{M_e,\prec}\) lies
in \(\mathrm{Non\text{-}sat}\), \(\mathrm{Fin\text{-}sat}\) or
\(\mathrm{Inf\text{-}axioms}\) according as \(e\in H_1\), \(e\in H_2\) or
\(e\in H_\infty\).  Lemma~\ref{thm:p02-w02-p02-036-gadget}(3)--(4) and
Lemma~\ref{thm:p02-w02-p02-036-transfer} preserve satisfiability and finite
satisfiability in both directions, so they preserve each of the three
classes.  Hence
\[
\begin{gathered}
 e\in H_1\Rightarrow g(e)\in\mathrm{Non\text{-}sat}(X),\qquad
 e\in H_2\Rightarrow g(e)\in\mathrm{Fin\text{-}sat}(X),\\
 e\in H_\infty\Rightarrow g(e)\in\mathrm{Inf\text{-}axioms}(X).
\end{gathered}
\]

\emph{Inseparability for \(X\).}  Pair each class with a halting set:
\(\mathrm{Non\text{-}sat}(X)\) with \(H_1\), \(\mathrm{Fin\text{-}sat}(X)\)
with \(H_2\), and \(\mathrm{Inf\text{-}axioms}(X)\) with \(H_\infty\).  Let
\(\mathcal C\) and \(\mathcal C'\) be two different classes, with halting sets
\(H\) and \(H'\), and suppose a decidable set \(W\) separates \(\mathcal C\)
from \(\mathcal C'\).  Put \(W'\defeq\{e\in V:g(e)\in W\}\); it is decidable
because \(V\) is decidable and \(g\) is computable on \(V\).  If \(e\in H\),
then \(g(e)\in\mathcal C\subseteq W\), so \(e\in W'\).  If \(e\in H'\), then
\(g(e)\in\mathcal C'\), which is disjoint from \(W\), so \(e\notin W'\).  So
\(W'\) separates \(H\) from \(H'\), contradicting
Lemma~\ref{thm:p02-w02-p02-036-halting}.

\emph{All formulas.}  Let \(\mathrm{FO}\) be the decidable set of codes of all
first-order formulas; by the coding fixed in the Inputs paragraph, \(X\) is a
decidable subset of it.  By the remark after
Definition~\ref{def:p02-w02-p02-036-classes},
\(\mathrm{Fin\text{-}sat}(X)=\mathrm{Fin\text{-}sat}\cap X\),
\(\mathrm{Non\text{-}sat}(X)=\mathrm{Non\text{-}sat}\cap X\) and
\(\mathrm{Inf\text{-}axioms}(X)=\mathrm{Inf\text{-}axioms}\cap X\); for a
sentence, satisfaction does not depend on the assignment
(\cref{lem:fol-coincidence}), so the formula notions agree with the sentence
notions.  If a decidable \(W\) separated one of the three classes of
\(\mathrm{FO}\) from another, then the decidable set \(W\cap X\) would contain
the corresponding class of \(X\) and be disjoint from the other, which was
just excluded.
\end{proof}
\end{annexitem}

\begin{annexitem}{annex:p02-w02-p02-036-exercise-2-1-29}{thm:p02-w02-p02-036-exercise-2-1-29}{Infinity axioms and unsatisfiable formulas}
\begin{proof}
The statement is the case \(\{\mathrm{Inf\text{-}axioms},
\mathrm{Non\text{-}sat}\}\) of Theorem~\ref{thm:p02-w02-p02-036-inseparable}.
The source's route, through formulas whose models are all infinite, is as
follows.  Let \(X=\operatorname{Sent}(\Lang)\) and \(R\) be as in that theorem,
and put \(g^{\infty}(e)\defeq(\rho^{\infty}_{M_e})^{\sharp}[R/E]\) for
\(e\in V\), a computable map by
Theorem~\ref{thm:p02-w02-p02-036-refined}, Lemma~\ref{thm:p02-w02-p02-036-gadget}
and the renaming.  By Theorem~\ref{thm:p02-w02-p02-036-refined}(4) and the
class preservation used in the proof of
Theorem~\ref{thm:p02-w02-p02-036-inseparable}, \(g^{\infty}(e)\in
\mathrm{Non\text{-}sat}(X)\) for \(e\in H_1\), and \(g^{\infty}(e)\in
\mathrm{Inf\text{-}axioms}(X)\) for \(e\in V\setminus H_1\).

Suppose a decidable \(W\) separates \(\mathrm{Inf\text{-}axioms}(X)\) from
\(\mathrm{Non\text{-}sat}(X)\), and put \(W'\defeq\{e\in V:g^{\infty}(e)\in
W\}\), a decidable set.  Then \(W'\) contains \(V\setminus H_1\) and is
disjoint from \(H_1\), so \(W'=V\setminus H_1\) and \(H_1=V\setminus W'\) is
decidable.  But then \(H_1\) itself is a decidable set separating \(H_1\) from
\(H_2\), contradicting Lemma~\ref{thm:p02-w02-p02-036-halting}.  The other
order follows by the symmetry remark after
Definition~\ref{def:p02-w02-p02-036-inseparable}, and the statement for all
formulas follows by intersecting a separating set with \(X\), as in the last
paragraph of the proof of Theorem~\ref{thm:p02-w02-p02-036-inseparable}.
\end{proof}
\end{annexitem}

\begin{annexitem}{annex:p02-w02-p02-036-conservative}{thm:p02-w02-p02-036-conservative}{Conservative reduction classes inherit inseparability}
\begin{proof}
Let \(g\) be a conservative reduction from \(\mathrm{FO}\) to \(X\).  For
\(\psi\in\mathrm{FO}\), the class of \(\psi\) among \(\mathrm{Fin\text{-}sat}\),
\(\mathrm{Inf\text{-}axioms}\) and \(\mathrm{Non\text{-}sat}\) is determined by
whether \(\psi\) is satisfiable and whether it is finitely satisfiable, and
\(g\) preserves both properties in both directions.  So
\(g(\mathrm{Fin\text{-}sat})\subseteq\mathrm{Fin\text{-}sat}(X)\),
\(g(\mathrm{Inf\text{-}axioms})\subseteq\mathrm{Inf\text{-}axioms}(X)\) and
\(g(\mathrm{Non\text{-}sat})\subseteq\mathrm{Non\text{-}sat}(X)\).

Let \(\mathcal C(X)\) and \(\mathcal C'(X)\) be two different classes of \(X\),
with \(\mathcal C\) and \(\mathcal C'\) the corresponding classes of
\(\mathrm{FO}\), and suppose a decidable set \(W\) separates \(\mathcal C(X)\)
from \(\mathcal C'(X)\).  Put
\(W'\defeq\{n\in\mathrm{FO}:g(n)\in W\}\).  It is decidable: decide
\(n\in\mathrm{FO}\), then compute \(g(n)\), which halts because
\(n\in\mathrm{FO}\), then decide \(g(n)\in W\).  If \(\psi\in\mathcal C\), then
\(g(\psi)\in\mathcal C(X)\subseteq W\), so \(\psi\in W'\).  If
\(\psi\in\mathcal C'\), then \(g(\psi)\in\mathcal C'(X)\), which is disjoint
from \(W\), so \(\psi\notin W'\).  So \(W'\) separates \(\mathcal C\) from
\(\mathcal C'\), contradicting Theorem~\ref{thm:p02-w02-p02-036-inseparable}.
\end{proof}
\end{annexitem}
